\documentclass[aspectratio=169,11pt]{beamer}
\usepackage{../../../shared/theme/esmad_beamer_theme}
\usepackage{amsmath,amssymb}
\usepackage{booktabs}

\title[Double Machine Learning]{Double Machine Learning and Orthogonal Scores}
\subtitle{Flexible Nuisance Estimation with Valid Causal Inference}
\date{\today}

\newcommand{\E}{\mathbb{E}}

\begin{document}

\begin{frame}
\titlepage
\end{frame}

\begin{frame}{Learning Goals}
\begin{itemize}
\item explain why naive plug-in machine learning can bias causal parameters;
\item derive the partialling-out idea in partially linear models;
\item explain Neyman orthogonality;
\item understand cross-fitting and why it helps;
\item distinguish nuisance estimation from target-parameter estimation;
\item recognize the identification assumptions that DML does not solve.
\end{itemize}
\end{frame}

\begin{frame}{Outline}
\tableofcontents
\end{frame}

\section{Why Double Machine Learning?}

\begin{frame}{Prediction Is Not the Target}
Suppose we want a low-dimensional causal parameter \(\theta_0\), but must also estimate flexible nuisance functions.

Examples:
\begin{itemize}
\item outcome regression;
\item propensity score;
\item conditional treatment model;
\item high-dimensional controls.
\end{itemize}

\begin{alertblock}{Problem}
Regularized ML can predict well while introducing bias into a plug-in causal estimator.
\end{alertblock}
\end{frame}

\begin{frame}{Partially Linear Model}
Consider
\[
Y = \theta_0 D + g_0(X) + \varepsilon,
\]
with
\[
\E[\varepsilon\mid X,D]=0.
\]

Treatment follows
\[
D = m_0(X)+v,
\]
with
\[
\E[v\mid X]=0.
\]
\end{frame}

\section{Partialling Out}

\begin{frame}{Residualize the Outcome}
Estimate
\[
g_0(X) \approx \widehat g(X),
\]
then form
\[
\widetilde Y = Y-\widehat g(X).
\]
\end{frame}

\begin{frame}{Residualize the Treatment}
Estimate
\[
m_0(X) \approx \widehat m(X),
\]
then form
\[
\widetilde D = D-\widehat m(X).
\]

Then estimate
\[
\widetilde Y = \theta_0 \widetilde D + \text{error}.
\]
\end{frame}

\begin{frame}{Why Partialling Out Helps}
The residualized treatment isolates variation in \(D\) that is not explained by \(X\).

The target parameter is then estimated from residualized outcome and residualized treatment rather than directly from a high-dimensional joint model.
\end{frame}

\section{Orthogonal Scores}

\begin{frame}{Moment Conditions}
Let the score be
\[
\psi(W;\theta,\eta),
\]
where \(\theta\) is the target and \(\eta\) contains nuisance functions.

The target solves
\[
\E[\psi(W;\theta_0,\eta_0)]=0.
\]
\end{frame}

\begin{frame}{Neyman Orthogonality}
An orthogonal score satisfies, informally,
\[
\left.
\frac{\partial}{\partial r}
\E[
\psi(W;\theta_0,\eta_0+r h)
]
\right|_{r=0}
=0.
\]

\begin{block}{Meaning}
Small first-order errors in nuisance estimation have little first-order effect on the estimating equation for \(\theta_0\).
\end{block}
\end{frame}

\begin{frame}{Orthogonality Is Not Magic}
Orthogonality reduces sensitivity to nuisance error.

It does not eliminate:
\begin{itemize}
\item unmeasured confounding;
\item positivity violations;
\item invalid instruments;
\item bad target definitions;
\item severe model misspecification outside the orthogonal structure.
\end{itemize}
\end{frame}

\section{Cross-Fitting}

\begin{frame}{Why Sample Splitting?}
Using the same sample to:
\begin{itemize}
\item fit complex nuisance models;
\item compute residuals;
\item estimate the target parameter
\end{itemize}
can produce overfitting bias.
\end{frame}

\begin{frame}{Cross-Fitting Procedure}
Split data into folds.

For each fold:
\begin{enumerate}
\item fit nuisance models on the other folds;
\item predict nuisance quantities on the held-out fold;
\item compute orthogonal scores on held-out observations.
\end{enumerate}

Finally aggregate scores across folds.
\end{frame}

\begin{frame}{Why Cross-Fitting Helps}
Cross-fitting weakens empirical-process restrictions and reduces dependence between nuisance fitting errors and score evaluation.

This permits flexible learners while retaining asymptotic inference under suitable rate conditions.
\end{frame}

\section{Rate Conditions}

\begin{frame}{Nuisance Models Need Not Be Parametric}
DML does not require perfect nuisance models.

A typical condition is that nuisance errors shrink sufficiently fast so their product is small enough, for example schematically:
\[
\|\widehat g-g_0\|
\cdot
\|\widehat m-m_0\|
=
o_p(n^{-1/2}).
\]
\end{frame}

\begin{frame}{Product-Rate Intuition}
One nuisance model can converge more slowly if the other is estimated very well.

\begin{alertblock}{But not arbitrarily slowly}
Orthogonality protects against first-order nuisance error, not against completely uninformative nuisance estimation.
\end{alertblock}
\end{frame}

\section{ATE Scores}

\begin{frame}{Propensity and Outcome Nuisance Functions}
For binary treatment:
\[
e(X)=P(D=1\mid X),
\]
\[
\mu_d(X)=\E[Y\mid D=d,X].
\]
\end{frame}

\begin{frame}{Orthogonal ATE Score}
An augmented score is
\[
\psi(W)
=
\mu_1(X)-\mu_0(X)
+
\frac{D(Y-\mu_1(X))}{e(X)}
-
\frac{(1-D)(Y-\mu_0(X))}{1-e(X)}.
\]

Its expectation equals the ATE under identification assumptions.
\end{frame}

\begin{frame}{Connection to Double Robustness}
This score combines:
\begin{itemize}
\item outcome regression;
\item propensity weighting.
\end{itemize}

It is doubly robust in the usual sense and orthogonal to small nuisance perturbations under regularity conditions.
\end{frame}

\section{Regularization Bias}

\begin{frame}{Why Lasso or Forest Plug-In Can Fail}
If we estimate a causal coefficient directly after aggressive regularization, shrinkage and selection can bias the target.

Prediction-optimal tuning is not necessarily inference-optimal.
\end{frame}

\begin{frame}{DML Strategy}
Separate:
\begin{itemize}
\item nuisance prediction;
\item orthogonal score construction;
\item low-dimensional target estimation.
\end{itemize}

This isolates the causal target from much of the regularization bias in nuisance models.
\end{frame}

\section{Inference}

\begin{frame}{Root-n Inference}
Under identification, orthogonality, overlap, and suitable nuisance-rate conditions,
\[
\sqrt{n}(\widehat\theta-\theta_0)
\]
can have an asymptotically normal limit.

This enables confidence intervals even with flexible nuisance learners.
\end{frame}

\begin{frame}{Clustered or Dependent Data}
Cross-fitting alone does not solve dependence.

Panel, grouped, or time-series data may require:
\begin{itemize}
\item clustered sample splitting;
\item cluster-robust variance estimation;
\item block-based cross-fitting.
\end{itemize}
\end{frame}

\section{Learner Choice}

\begin{frame}{What Can Estimate the Nuisance Functions?}
Possible learners include:
\begin{itemize}
\item lasso / elastic net;
\item random forests;
\item boosted trees;
\item neural networks;
\item splines and GAMs;
\item ensembles.
\end{itemize}

The causal target remains defined independently of the learner.
\end{frame}

\begin{frame}{Prediction Quality Still Matters}
Orthogonality is not an excuse for weak nuisance models.

Poor nuisance estimation can cause:
\begin{itemize}
\item unstable residualization;
\item overlap problems;
\item poor finite-sample performance;
\item invalid asymptotic approximations.
\end{itemize}
\end{frame}

\section{Diagnostics}

\begin{frame}{Overlap Diagnostics}
For treatment-effect DML, inspect the estimated propensity distribution.

Extreme values near 0 or 1 produce unstable scores and large variance.
\end{frame}

\begin{frame}{Score Diagnostics}
Useful checks include:
\begin{itemize}
\item fold-specific estimates;
\item nuisance prediction errors;
\item score outliers;
\item sensitivity to learner choice;
\item sensitivity to cross-fitting partitions.
\end{itemize}
\end{frame}

\section{Applied Reasoning}

\begin{frame}{Example: Marketing Treatment Effect}
Suppose treatment assignment depends on hundreds of behavioral features.

A DML workflow:
\begin{itemize}
\item fit flexible propensity and outcome models;
\item cross-fit predictions;
\item compute orthogonal ATE or ATT scores;
\item estimate the low-dimensional treatment effect;
\item report uncertainty and overlap diagnostics.
\end{itemize}
\end{frame}

\begin{frame}{A Credible DML Workflow}
\begin{enumerate}
\item Define the causal estimand and identification assumptions.
\item Choose an orthogonal score for the target.
\item Specify nuisance functions.
\item Fit nuisance learners using training folds only.
\item Cross-fit predictions to held-out folds.
\item Compute orthogonal scores.
\item Estimate the target parameter and variance.
\item Inspect overlap and score stability.
\item Repeat across defensible nuisance learners.
\end{enumerate}
\end{frame}

\begin{frame}{Common Failure Modes}
\begin{itemize}
\item believing ML automatically removes confounding;
\item using flexible models without cross-fitting;
\item treating orthogonality as immunity to bad overlap;
\item ignoring dependence structure;
\item reporting only one nuisance learner;
\item confusing predictive performance with causal identification.
\end{itemize}
\end{frame}

\begin{frame}{Synthesis: ML for Nuisance, Causality for the Target}
\begin{itemize}
\item DML separates nuisance prediction from target estimation.
\item Partialling out removes predictable components of treatment and outcome.
\item Orthogonal scores reduce first-order sensitivity to nuisance errors.
\item Cross-fitting prevents the same observations from driving nuisance fitting and score evaluation.
\item Flexible ML can be used while retaining classical inference under suitable conditions.
\item DML does not replace identification assumptions.
\item Overlap, dependence, and target definition remain central.
\end{itemize}
\end{frame}

\begin{frame}{Selected References}
\small
\begin{itemize}
\item Chernozhukov, V. et al. (2018). Double/Debiased Machine Learning for Treatment and Structural Parameters. \textit{The Econometrics Journal}, 21(1), C1--C68.
\item Robinson, P. M. (1988). Root-N-Consistent Semiparametric Regression. \textit{Econometrica}, 56(4), 931--954.
\item Belloni, A., Chernozhukov, V., and Hansen, C. (2014). Inference on Treatment Effects after Selection among High-Dimensional Controls. \textit{Review of Economic Studies}, 81(2), 608--650.
\item Chernozhukov, V. et al. (2022). Locally Robust Semiparametric Estimation. \textit{Econometrica}, 90(4), 1501--1535.
\end{itemize}
\end{frame}

\contactslide

\end{document}
